\documentclass[11pt,a4paper]{article}
\usepackage[T1]{fontenc}
\usepackage[utf8]{inputenc}
\usepackage{lmodern,amsmath,amssymb,amsthm,mathtools,microtype,mathrsfs}
\usepackage[margin=25mm,top=23mm,bottom=25mm]{geometry}
\usepackage{booktabs,array,longtable,enumitem,xcolor,fancyhdr}
\usepackage[colorlinks=true,linkcolor=blue!45!black,urlcolor=blue!45!black,citecolor=blue!45!black]{hyperref}
\hypersetup{pdftitle={Reading the Alpoge--Claude RH and Dirichlet-GRH proof: working version 0.44},pdfauthor={ChatGPT},pdfsubject={Partial human-readable reconstruction with explicit arithmetic inputs},pdfkeywords={Riemann hypothesis, Lean, theta lifts, prime sums, Arakelov slopes}}
\setlength{\parindent}{0pt}
\setlength{\parskip}{5pt plus 1pt}
\setlist{topsep=3pt,itemsep=3pt,leftmargin=*}
\setlength{\headheight}{14pt}
\pagestyle{fancy}\fancyhf{}
\fancyhead[L]{\small RH and Dirichlet GRH: a reconstruction}
\fancyhead[R]{\small Working version 0.44}
\fancyfoot[C]{\thepage}
\renewcommand{\headrulewidth}{0.3pt}
\numberwithin{equation}{section}
\newtheorem{theorem}{Theorem}[section]
\newtheorem{proposition}[theorem]{Proposition}
\newtheorem{lemma}[theorem]{Lemma}
\newtheorem{corollary}[theorem]{Corollary}
\theoremstyle{remark}\newtheorem{remark}[theorem]{Remark}
\newcommand{\Rea}{\operatorname{Re}}
\newcommand{\HS}{\mathrm{HS}}
\newcommand{\tr}{\operatorname{tr}}
\newcommand{\codim}{\operatorname{codim}}
\newcommand{\covol}{\operatorname{covol}}
\newcommand{\OK}{\mathcal O_K}
\newcommand{\src}[2]{\href{https://github.com/CocoSgt/riemann-hypothesis-lean/blob/7cbf7a0c32fca1badf05931879a37c3f8a2b03b4/src/#1}{\texttt{#2}}}
\newcommand{\status}[1]{\textbf{Status:} #1}

\begin{document}
\section{One actual ordinary module family at every auxiliary level}
\label{sec:hida-level-modules-v45}

\status{HLM1--HLM26 reconstruct the full original \texttt{LevelAlt}
producer.  The same Taylor--Wiles base is chosen before all auxiliary
levels, including the empty auxiliary set.  Every module vector is
separated by the original classical points, with their full local
correspondence quantifiers.  The original \(\mathrm{HypothesesH}'\)
and its named local, automorphic, reciprocity and cohomological
foundations remain hypotheses.  No formal theorem is substituted
for the arguments below.}

\subsection{Notation and retained inputs}

Use the original field \(L\), closure \(\Omega\), setting \(s\),
coefficient prime \(\ell\), coefficient DVR \(O=c.O\), rank
\(n\geq1\), and fixed algebraic isomorphism
\(\iota:\overline{\mathbb Q}_\ell\simeq\mathbb C\).
There is no continuity assumption on \(\iota\), no restriction
\(\ell>n\), and no exclusion of \(\ell=3\).
For an original Hecke setting \(H\), a complete Noetherian local
\(O\)-algebra \(A\) in \(\mathcal C_O\), and the given weight
and residual eigenvalues \(\theta,a\), put
\[
 N_K=H.X.\operatorname{Ord}(H.\tau_A,\theta,H.\operatorname{lev}K),
 \qquad M_K=H.\operatorname{MK}(\theta,a,K),
 \qquad T_K=H.\operatorname{TK}(\theta,a,K).
 \tag{HLM.1}
\]
These are the original submodules of one coefficient-valued
function space, namely \(F=H.X.\operatorname{Fn}(A,H.m)\).
An auxiliary level is a tuple of compact open subgroups
\(K_q\subset G_q=\operatorname{GL}_n(L_{H.wq(q)})\);
write \(K=\prod_qK_q\) when no confusion results.
The residual part of a finite stable module \(N\), for commuting
operators \(T_i\) and scalars \(a_i\), means exactly
\[
 \operatorname{loc}_a(N)=
 \{z\in N:\ \forall i,k\ \exists r\geq0,
                    (T_i-a_i)^rz\in\mathfrak m_A^kN\}.
 \tag{HLM.2}
\]
Thus \(M_K=\operatorname{loc}_a(N_K)\), with the scalar
\(a_i\) mapped from \(O\) to \(A\).

OST19, OST24 and OHA5--OHA11 supply the actual ordinary projector,
finite freeness, arbitrary coefficient base change and residual
idempotents.  ASP13 supplies one specialization map on all levels.
AFS5, AFS19--AFS20 and UPA17 supply the density, part action and
actual automorphic parts when the setting comes from
\(\mathrm{UnitaryHeckeData}\).  LTP supplies the original sections
and their exact local actions.  Their hypotheses are retained,
including the original finite-level and neatness assertions FNL16
and NTL17.  The abstract statements below print any such hypotheses
when they are not yet installed.

The use of local class field theory is its original full
\(\mathrm{LocCFT}\)/Artin interface.  Existence of the auxiliary
universal determinant ring uses the original cohomological
\(\Phi\)-finiteness package, already reconstructed in KX6--KX9
from its stated inputs.  Neither a new local correspondence nor a
new cohomological finiteness theorem is asserted here.

\subsection{Projectivity and the actual directed union}

% BEGIN HLM1
\paragraph{HLM1. An integral trace criterion for projectivity.}
\label{claim:HLM1}
Let a finite group \(\Delta\) act on an \(A\)-projective module
\(N\) by \(\rho\).  If an \(A\)-linear map \(u:N\to N\)
satisfies
\[
             \sum_{g\in\Delta}\rho(g)u\rho(g^{-1})=1,
 \tag{HLM.3}
\]
then \(N\) is projective over \(A[\Delta]\).
Its norm map has image \(N^\Delta\) and kernel generated over
\(A\) by the differences \(\rho(g)z-z\).
\begin{proof}
The induced module \(A[\Delta]\otimes_A N\), with action on
the first factor, is projective: a splitting of an \(A\)-free
module onto \(N\) induces a splitting of free group-algebra
modules.  The equivariant surjection
\(g\otimes z\mapsto\rho(g)z\) has equivariant section
\[
               z\longmapsto
                 \sum_g g\otimes u(\rho(g^{-1})z).
\]
Equivariance follows by replacing \(g\) by \(hg\), and the
composite is (HLM.3).  This proves projectivity.
If \(z\) is invariant, (HLM.3) reads
\(z=\sum_g\rho(g)u(z)\), proving surjectivity of the norm
onto the invariants.  If \(\sum_g\rho(g)z=0\), then
\(\sum_g u(\rho(g^{-1})z)=0\); subtract this from (HLM.3)
applied to \(z\).  The result expresses \(z\) as a sum of
differences.  Conversely every difference is killed by the norm,
by reindexing its sum.  No division by \(|\Delta|\) occurs.
\end{proof}
% END HLM1

% BEGIN HLM2
\paragraph{HLM2. A neat fundamental domain gives the required trace.}
\label{claim:HLM2}
For the original frame \(X\), let \(V\triangleleft U\) be
finite-index subgroups of \(X.U_b\), and assume \(X.\mathrm{Neat}(U)\).
There is a function \(\chi:G_a\times G_p\to\{0,1\}\subset A\)
which is invariant under left \(\Gamma\), right \(V\), and
right \(I(0)\), and satisfies
\[
     \sum_{q\in U/V}\chi(x(q,1))=1\quad\hbox{for every }x.
 \tag{HLM.4}
\]
Multiplication by \(\chi\) preserves the original space
\(\operatorname{Sp}(\tau,\theta,V)\).
\begin{proof}
Choose one representative in each equivalence class for
\(y=X.\mathrm{glob}(\gamma)x(u,b)\),
\(u\in U,b\in I(0)\).  If two presentations from the same
chosen representative have respective \(u,u'\), cancellation
gives a stabilizer relation with away component \(u'u^{-1}\).
Neatness says that component is one.  Thus the away component
\(u\), and in particular its coset modulo \(V\), is unique.
Let \(\chi\) be one precisely when this coset is trivial.
Left global translation and right \(I(0)\) do not change it;
right \(V\) does not change whether it is trivial.  For
\(xq\), its coset is the coset of \(uq\), so precisely one
\(q\in U/V\) contributes to (HLM.4).

The defining equations of \(\operatorname{Sp}\) are global
invariance, invariance under the positive unipotent, the specified
torus character, congruence invariance under a sufficiently deep
opposite unipotent at each precision, and the full matrix-valued
\(V\)-equation.  Every subgroup occurring in the first four
right equations is contained in \(I(0)\).  Hence multiplication
by \(\chi\) commutes with those translations and preserves
their congruences.  In the last equation the scalar \(\chi\)
commutes with the coefficient matrix \(\tau(u)\).  This checks
all equations, including the original coefficient action.
\end{proof}
% END HLM2

% BEGIN HLM3
\paragraph{HLM3. The full ordinary finite-quotient projectivity argument.}
\label{claim:HLM3}
Assume the original finite-level, continuous-weight and neatness
hypotheses at \(U\), and \(V\triangleleft U\) of finite index.
Then \(\operatorname{Ord}(\tau,\theta,V)\) is projective
over \(A[U/V]\) for the actual right-translation action.
The original double-coset trace has image
\(\operatorname{Ord}(\tau,\theta,U)\) and kernel the span
of the translation differences at level \(V\).
\begin{proof}
Finite index transfers the finite class bound to \(V\);
neatness transfers the stabilizer freeness condition.  OST19
therefore gives finite \(A\)-freeness and the ordinary projector
\(e:\operatorname{Sp}(V)\to\operatorname{Ord}(V)\).
Right translations by \(U\) preserve \(\operatorname{Sp}(V)\)
and its ordinary subspace, because \(V\) is normal.  They
commute with \(U_z\), and preserve every value-congruence
filtration.  The factorial-limit characterization of \(e\)
in OST12 consequently gives \(e\rho(u)=\rho(u)e\): compare
the finite powers at each precision and then use separatedness.
No topology on the coefficient isomorphism is involved.

Let \(m_\chi\) be HLM2's multiplication map and put
\(u=e m_\chi\) on \(\operatorname{Ord}(V)\).
Before applying \(e\), (HLM.4) gives, pointwise and with the
matrix factors cancelling,
\(\sum_q\rho(q)m_\chi\rho(q^{-1})=1\).
Commuting \(e\) through this finite sum gives (HLM.3).
HLM1 proves projectivity and the two norm assertions.
Finally the original unnormalized double-coset operator at one
is exactly \(\sum_{q\in U/V}\rho(q)\) on \(V\)-fixed
vectors: its chosen representatives enumerate those right
cosets, and changing a representative multiplies by an element
fixing the vector.  The \(U\)-invariants of \(\operatorname{Ord}(V)\)
are \(\operatorname{Ord}(U)\), since their ambient and
ordinary congruences are identical.  This proves all parts of
the original \texttt{ordProj} statement at its stated scope.
\end{proof}
% END HLM3

% BEGIN HLM4
\paragraph{HLM4. Residual parts commute with injections and intersections.}
\label{claim:HLM4}
Let \(N\subset N'\) be stable finite modules over a complete
Noetherian local ring, with the same commuting operators.
Then
\[
       \operatorname{loc}_a(N)=N\cap\operatorname{loc}_a(N').
 \tag{HLM.5}
\]
An operator-intertwining linear map carries residual parts into
residual parts.  The canonical residual projector commutes with
every endomorphism commuting with the operators.
\begin{proof}
Functoriality follows directly from (HLM.2).  For the reverse
inclusion in (HLM.5), the Artin--Rees argument gives a fixed
\(c\) such that
\(N\cap\mathfrak m^{k+c}N'\subseteq\mathfrak m^kN\)
for all \(k\).  Here is the needed finite-generation argument:
the Rees ring \(\bigoplus_k\mathfrak m^kt^k\) is Noetherian,
being generated by the finitely many degree-one generators of
\(\mathfrak m\); its finite module
\(\bigoplus_k\mathfrak m^kN't^k\) has finite graded submodule
\(\bigoplus_k(N\cap\mathfrak m^kN')t^k\).
Take \(c\) at least its generator degrees.  Multiplying those
generators proves the displayed inclusion.  If \(z\in N\)
satisfies (HLM.2) in \(N'\), apply its condition at precision
\(k+c\); the relevant operator power is also in \(N\),
so it is in \(\mathfrak m^kN\).
Finally OHA5--OHA6 construct the residual projector as an
idempotent of the finite algebra generated by the operators.
Every commuting endomorphism commutes with this idempotent.
\end{proof}
% END HLM4

% BEGIN HLM5
\paragraph{HLM5. The actual smooth union and all its fixed modules.}
\label{claim:HLM5}
Under the original \(H.\mathrm{Ax}\) and \(\mathrm{IsWt}\),
put \(\mathcal M=\bigcup_KM_K\subset F\).
It is a module with a representation \(\rho\) of
\(G_Q=\prod_qG_q\), characterized by
\[
 \rho(g)f=H.X.\mathrm{ra}(H.\tau_A)(j(g))f,
 \qquad j(g)=\prod_qH.\mathrm{sec}(H.wq(q))(g_q).
 \tag{HLM.6}
\]
Every finite set of its vectors lies in one \(M_K\), and
\(\mathcal M^K=M_K\) as original functions.  In particular
these fixed modules are finite over \(A\).
If all \(H.\mathrm{lev}K\) are neat and the original section
faithfulness condition holds, then for every \(K'\triangleleft K\)
the module \(\mathcal M^{K'}\) is projective over
\(A[K/K']\), with the quotient action induced by \(\rho\).
\begin{proof}
The section images at different places commute by \(H.\mathrm{Ax}\).
Thus the finite product \(j\) is a homomorphism and
\[
       H.\mathrm{lev}K=H.U_Q\,j(K),
              \qquad [H.U_Q,j(G_Q)]=1.
 \tag{HLM.7}
\]
Intersections and conjugates of compact open level tuples are
again such tuples; their indices in a compact level are finite
because the corresponding compact quotient is discrete.
If \(K'\subset K\), then \(N_K\subset N_{K'}\), and
HLM4 gives \(M_K\subset M_{K'}\).  Intersections of levels
therefore prove directedness and the finite-set assertion.

Right translation by \(j(g)\) takes \(N_K\) to
\(N_{gKg^{-1}}\).  This follows from (HLM.7), the original
ambient equations, and commutation with \(U_z\).  It also
commutes with every original Hecke operator: away sections
commute at different places, while the \(\ell\)-adic
operators act on the other factor of the function argument.
HLM4 gives stability of the residual parts, proving (HLM.6).
Its order is the source order
\(\rho(g)\rho(h)=\rho(gh)\), inherited from
\(\mathrm{ra}(u)f(x)=\tau(u)f(x(u,1))\).

Certainly \(M_K\subset\mathcal M^K\).  Conversely take
\(f\in M_{K_1}\cap\mathcal M^K\) and refine both levels
to \(K_2\).  The vector lies in \(N_K\), since it is fixed
by both \(U_Q\) and \(j(K)\); all other defining conditions
are already those of \(N_{K_2}\).
HLM4, applied to \(N_K\subset N_{K_2}\), now gives
\(f\in M_K\).  OHA8 gives the asserted finiteness.

For projectivity, (HLM.7) and faithfulness identify
\(K/K'\) with
\(H.\mathrm{lev}K/H.\mathrm{lev}K'\): surjectivity removes
the common \(U_Q\), and the kernel is exactly \(K'\).
Normality follows from commutation in (HLM.7).
HLM3 makes \(N_{K'}\) projective over this group algebra.
Its residual projector commutes with \(K/K'\) by HLM4.
Consequently \(M_{K'}=\mathcal M^{K'}\) is an equivariant
direct summand and is projective.  This proves the complete
\texttt{piBasic} output, including the actual fixed-module
identification, rather than just existence of some smooth module.
\end{proof}
% END HLM5

\subsection{Universal Hecke maps at every level}

% BEGIN HLM6
\paragraph{HLM6. Integral determinant interpolation in a compact coefficient ring.}
\label{claim:HLM6}
Let \(T\in\mathcal C_O\), and let a family
\(x:T\to_O\overline{\mathbb Q}_\ell\) separate its elements.
Suppose \(s_j\) is dense in a topological group \(\Gamma\),
and polynomials \(P_j\in T[X]\) specialize, for each \(x\),
to the characteristic polynomials of one continuous
\(n\)-dimensional representation \(r_x(s_j)\).
There is a continuous degree-\(n\) determinant law \(D\) over
\(T\) with \(D\otimes x=\det r_x\) and
\(\chi_D(s_j)=P_j\).
Two determinant laws over any ring are equal if all their
characteristic coefficients on group elements are equal.
\begin{proof}
We recall the integral word identity used in JSP9.  For finitely
many noncommuting letters, the unique ordered Lyndon factorization
gives a formal factorization of \(1-\sum Z_i[g_i]\) into
factors \(1-Z_w[g_w]\) indexed by Lyndon words.  This identity
can be checked degree by degree: at each degree, the leading
remaining word determines the next factor, and unique Lyndon
factorization gives coefficient one for every word of the inverse
series.  Only finitely many factors contribute in any fixed
degree.  Work in each finite-degree truncated scalar algebra:
there the product and every group-algebra element to which the
determinant is applied have finite support.  Apply the
multiplicative homogeneous determinant law there and discard
terms above its degree.  No law on an infinite-support
group-algebra element is presumed.  The coefficients of its value on
\(1-\sum Z_i[g_i]\) are integral polynomials in the
characteristic coefficients of the group words \(g_w\).
The top homogeneous part gives the value on
\(\sum Z_i[g_i]\).  Naturality then permits arbitrary
coefficients in every scalar algebra.  This proves the last
assertion and supplies the required integral formulas.  In
particular, no Newton division in \(T\) is used.

The maximal-ideal quotients of \(T\) are finite: its residue
field is finite, and each successive ideal quotient is finite
dimensional over it.  Completeness and separatedness identify
\(T\) with their inverse limit, so its adic topology is compact
Hausdorff.  ASP2 shows that each given \(x\) is continuous;
indeed it bounds the image of \(\mathfrak m_T^k\) by
\(b_x^k\) with \(b_x<1\).  Thus
\(T\hookrightarrow\prod_x\overline{\mathbb Q}_\ell\)
is a homeomorphism onto a closed subring.
For each \(g\in\Gamma\), the tuple of the \(i\)th
characteristic coefficients of \(r_x(g)\) lies in that subring:
this is true at the dense \(s_j\), and the tuple varies
continuously.  Its unique inverse image \(c_i(g)\in T\)
is therefore continuous.

Use the integral word formulas just proved to define the universal
polynomial associated with every finite sum of group elements,
using these \(c_i(g_w)\).  After every \(x\), it is exactly
the polynomial for \(\det r_x\).  Separation applied to
polynomial coefficients shows that these polynomials are
consistent on overlapping lists, homogeneous of degree \(n\),
multiplicative and unital.  Substitution defines a natural law
on every commutative \(T\)-algebra.  This proves existence of
\(D\), its full scalar-extension identities and continuity.
The specified Frobenius polynomials are recovered coefficientwise
by the same separation.  This proof does not assert injectivity
after tensoring \(T\) into the product of fields.
\end{proof}
% END HLM6

% BEGIN HLM7
\paragraph{HLM7. The full original \texttt{coreDet} producer.}
\label{claim:HLM7}
Assume the original \(H.\mathrm{Ax}\) and \(\mathrm{IsWt}\),
the supplied automorphic-set hypothesis
\[
            H.\mathrm{AutomorphicHeckeSet}(\iota,P),
\]
an admitted \(\lambda_0\) with \(\mathrm{DenseWt}(\theta,\lambda_0)\),
and the original local Hecke algebra \(T_K\in\mathcal C_O\).
There are \(D_K,\psi_{K,v,i}\) satisfying every field of
\(H.\mathrm{IsUniv}\) on this very algebra.  This abstract
supplier also permits \(n=0\), although the final unitary
\(\mathrm{LevelAlt}\) producer assumes \(n\geq1\).
\begin{proof}
If \(n=0\), take the constant degree-zero determinant
law \(1\); its characteristic polynomial and the original
degree-zero Hecke polynomial are both \(1\), and the local
character family is empty.  All five fields follow directly.
We may therefore assume \(n\geq1\).
OHA10 makes every positive generator
\(U_{v,j}=\mathrm{gen}(\mathrm{inr}(v,j))\) a unit of
\(T_K\).  Put \(U_{v,-1}=1\) and
\(u_{v,i}=U_{v,i}U_{v,i-1}^{-1}\).
On the unit group of \(L_{H.S_l(v)}\), prescribe
\(\xi_{v,i}(y)=\theta(e_{v,i}(y))^{-1}\), now mapped
to \(T_K\).  Extend to the multiplicative group by
\[
 \xi_{v,i}(y\varpi^a)=\xi_{v,i}(y)u_{v,i}^a.
 \tag{HLM.8}
\]
The valuation decomposition makes this a character.  The original
\(\mathrm{ContChar}\) condition says its reduction modulo each
\(\mathfrak m_{T_K}^k\) is continuous.  That quotient is finite,
so the full local reciprocity interface extends it uniquely to
a continuous character of the local abelian Galois group.
Taking the inverse limit and multiplying by
\(\epsilon^{-i}\) defines \(\psi_{K,v,i}\) on the actual
decomposition group.  The unit formula is immediate, and the
uniformizer formula is the telescoping identity
\[
 \prod_{i\leq j}\psi_{K,v,i}(g)\epsilon(g)^i=U_{v,j}
      \quad\text{if }\operatorname{Art}(g)=\varpi.
 \tag{HLM.9}
\]
This proves both original character formulas and their adic
continuity, on every permitted Weil lift.

AFS20 says that the good \(T_K\)-points at \(\lambda_0\)
separate \(T_K\).  For each specified such point, ASP13 and
AFS19(3) give a nonzero eigenvector in an actual part.  Its
representation is continuous, and its spherical eigenvalues are
the point's values of the exact original generators.  The part's
Frobenius equation therefore says that its characteristic
polynomial at \(H.\mathrm{prj}(\sigma^{-1})\) is
\[
 \operatorname{heckePol}\bigl(n,\operatorname{Norm}w,
                               x(\operatorname{genW}(w,-))\bigr).
\]
The degree-zero term and indices \(j+1\) are those of the
original \(\operatorname{genW}\); no spherical normalization
is changed.  \(H.\mathrm{Ax.dense}\) and HLM6 produce
\(D_K\) with this polynomial and adic continuity.  Along
with (HLM.8)--(HLM.9), these are exactly the five
\(\mathrm{IsUniv}\) fields.  The construction precedes any
use of an \(\mathrm{Out}\) record that assumes such a pair.
\end{proof}
% END HLM7

% BEGIN HLM8
\paragraph{HLM8. Uniqueness and functoriality of the original universal action.}
\label{claim:HLM8}
Two \(\mathrm{IsUniv}\) pairs on the same original Hecke
algebra are equal.  If an injective \(A\)-linear map
\(f:M_K\to M_{K'}\) intertwines every original Hecke
generator, and a ring \(R\) with pair \((D_0,\psi_0)\)
has the source's unique-map property into \(\mathcal C_O\),
then any given \(\mathrm{IsMap}\) homomorphisms from \(R\)
to the two Hecke algebras induce intertwining actions on \(f\).
\begin{proof}
The two determinant characteristic coefficients agree on the
specified dense Frobenius set.  Modulo each maximal-ideal power
they are continuous functions to a finite discrete ring, so they
agree everywhere.  Separatedness and HLM6 give equality of full
determinant laws.  The unit formulas determine each local
character on units.  Equality of the prefix products in the
uniformizer formulas determines the individual values by
cancellation in the unit group.  Units and one uniformizer
generate the local multiplicative group, so reciprocity gives
equality on the Weil subgroup.  Its density, continuity modulo
every ideal power, and separatedness give equality on the
decomposition group.

For the action assertion, first construct
\(\Phi:T_{K'}\to_A T_K\).  Every polynomial in the target
operators preserves \(f(M_K)\); its restriction there is the
same polynomial in the source operators.  If two polynomials
give the same target operator, injectivity of \(f\) makes
their source operators equal.  This defines \(\Phi\) on the
algebra generated by the operators, preserving all scalars and
generators, and gives
\(f\Phi(t)=tf\).  It is an \(O\)-algebra map between
\(\mathcal C_O\) objects, hence local and adically continuous:
the common residue maps give
\(\Phi(\mathfrak m^k)\subset\mathfrak m^k\).
Pushing the target \(\mathrm{IsUniv}\) pair through \(\Phi\)
preserves all five fields, since it preserves the relevant
scalars, units and generators.  Uniqueness makes it the source
pair.  The unique-map property of \(R\) then gives
\(f_K=\Phi\circ f_{K'}\), which is precisely the asserted
intertwining.  If the source module is zero the assertion is
automatic; an injection from a nonzero source excludes a zero
target.  This proves the literal \texttt{TCoreUniq} and
\texttt{LevelUR2Compat} consumers without an extra density
hypothesis on \(A\).
\end{proof}
% END HLM8

% BEGIN HLM9
\paragraph{HLM9. Original universal Hecke maps and ideals exist at every level.}
\label{claim:HLM9}
For every original supplied \(\mathrm{UnitaryHeckeData}\),
universal ordinary determinant ring \(B\), weight object
\(\Lambda\), and level \(K\), there is an ideal satisfying
the full original \(\mathrm{IsJ}\).  If \(M_K\ne0\), its
map is a surjection \(f_K:B\to T_K\), is an
\(\mathrm{IsMap}\), and satisfies
\(f_K\Lambda.\mathrm{toPair}(B.u)=A\to T_K\).
\begin{proof}
If \(M_K=0\), choose \(J=B\), the literal zero branch.
Otherwise OHA11 gives the nonzero \(\mathcal C_O\) algebra
\(T_K\) with its exact residual eigenvalues.  FNL16,
UPA17 and AFS5 install HLM7 at the admitted zero weight.
It gives \(D_K,\psi_K\); HEP11 now applies to this supplied
pair and gives its actual separating good points and their
continuous representations and flags.

We check that this is an original ordinary determinant pair.
Modulo \(\mathfrak m_{T_K}\), the Frobenius formula and
\(\mathrm{IsRes.frob}\) agree.  Their coefficient functions
are locally constant on the actual group.  Density and HLM6
give \(\overline D_K=r.\overline D\) as full laws.
On inertia, compare the \(\mathrm{IsTheta}\) equation with
the unit equation (HLM.8); cancellation gives
\(\psi_{K,v,i}|_{I_v}=\kappa_{v,i}\) after the given
scalar map.  Reduction of (HLM.9) and the literal residual
uniformizer formula identify the residual characters on a
uniformizer.  Their agreement on inertia, then on the dense
Weil subgroup, gives the full prescribed reduction.  For the
torsion condition, the original \(\Lambda.\mathrm{inert}\)
or \(\mathrm{IsKap}\) field supplies a continuous extension
\(\psi^0_{v,i}\) of \(\kappa_{v,i}\) having the prescribed
residual and torsion values; map it to \(T_K\).  The quotient
\(\psi_{K,v,i}/\psi^0_{v,i}\) is trivial on inertia.
The actual local decomposition group modulo inertia is the
unramified group \(\widehat{\mathbb Z}\), which is
torsion-free.  An element of the source \(\mu_G\) has
\(\ell\)-power-torsion image in the local abelianization,
so its image in this unramified quotient is zero.  The
quotient character therefore kills the full preimage
\(\mu_G\), exactly as required by the original
\(\mathrm{IsPsi.mu}\) field.  Thus every \(\mathrm{IsPsi}\) field
holds, and the weight map is the original one.

The remaining pair relations can be checked at the separating
good points.  The part's form makes the determinant invariant
under the prescribed star.  The ordered flag kills the ordered
ordinary product, in its source order.  For each residual type
place choose the original Teichmuller lifts of the residual
characters.  The part's type diagonalization makes every
element of the source \(I_p(u)\) act trivially, because its
residual characters are all one and the Teichmuller lift of one
is one.  Thus the required trace relations hold at every point.
The universal coefficient argument of JRP6--JRP12, or directly
the polynomial-coefficient argument of JSP2, descends these
relations to \(T_K\).  This constructs all fields of the
original pair, not just its values at points.

Universality of \(B\) gives \(f_K\) carrying its pair to
\((D_K,\psi_K)\).  Universality of \(\Lambda\), using the
inertia identity already proved, gives weight compatibility.
Every \(U_{v,j}\) is the prefix product (HLM.9), hence lies
in the image.  The \(j\)th away generator is the
\((n-j-1)\)st characteristic coefficient multiplied by
\((-1)^{j+1}q^{-(j+1)j/2}\).  This scalar is an original
\(O\)-unit because that place is away from \(\ell\).
Consequently every Hecke generator and every \(A\)-scalar
is in the image.  They generate \(T_K\), so \(f_K\) is
surjective.  Taking \(J=\ker f_K\) proves the complete
nonzero \(\mathrm{IsJ}\) branch.  This argument is valid
for all original ranks, not just the rank-four application
of JSP3--JSP5.
\end{proof}
% END HLM9

% BEGIN HLM10
\paragraph{HLM10. One universal-ring action on the entire union.}
\label{claim:HLM10}
Assume the original abstract \texttt{umodAbs} hypotheses:
\((R,D_0,\psi_0)\) has the unique-map property of HLM8,
\(\lambda:A\to_O R\) is given, and for every \(K\)
there exists an original \(\mathrm{IsJ}\) for these data.
Then \(\mathcal M\) of HLM5 has one ring action
\(\alpha:R\to\operatorname{End}_A(\mathcal M)\), with
\[
 \alpha(\lambda(b))z=bz,\qquad
 \alpha(t)\rho(g)z=\rho(g)\alpha(t)z.
 \tag{HLM.10}
\]
Its restriction to each nonzero \(M_K\) is induced by an
original \(\mathrm{IsMap}\) homomorphism \(R\to T_K\).
\begin{proof}
Choose the action \(\mu_K\) furnished by the stated
\(\mathrm{IsJ}\), using the unique zero action if \(M_K=0\).
It has the first identity of (HLM.10).  For an inclusion
\(M_K\subset M_{K'}\), the operators are the same actual
operators on functions.  HLM8 says that the actions agree.
For \(z\in\mathcal M\), define \(\alpha(t)z=\mu_K(t)z\)
at any level containing it.  A common refinement proves
independence, additivity, \(A\)-linearity, and both ring laws;
all involve finitely many vectors, so one such refinement
suffices for each equality.

For a group element \(g\), translation is an injective
Hecke-intertwining map \(M_K\to M_{gKg^{-1}}\), with
inverse translation by \(g^{-1}\).  HLM8 again proves
the second identity of (HLM.10).  The restriction formula is
the definition, including its equality in the original ambient
function space.  Thus no level-dependent action is being
silently identified with another one.
\end{proof}
% END HLM10

\subsection{Classical points separate all module vectors}

% BEGIN HLM11
\paragraph{HLM11. One scalar-character projection on a directed union.}
\label{claim:HLM11}
Let an additive map \(s:M\to V\) take values in a complex vector
space, let \(M=\bigcup_KM_K\) be directed, and suppose the
actions \(\alpha(t)\) and \(\rho(g)\) commute.  Suppose
\(s\rho(g)=\pi(g)s\).  Put
\(E_K=\operatorname{span}_{\mathbb C}s(M_K)\).
Assume every \(E_K\) under consideration has an action of
\(R\) agreeing with \(s\alpha\), and is spanned by
subspaces contained in specified parts \(Y_i\), on each of
which this action is scalar through a character
\(\chi_x:R\to\mathbb C\) related to that part.
It suffices to have these assertions at all levels containing
one specified \(z\).  If \(s(z)\ne0\), there are a character
\(\chi_{x_0}\), a \(\pi\)-stable subspace \(E_0\), and
a complex linear projection on \(E=\operatorname{span}s(M)\)
onto \(E_0\), such that
\[
 p(s(z))\ne0,\qquad p(s(\alpha(t)m))=\chi_{x_0}(t)p(s(m)),
 \qquad p\pi(g)=\pi(g)p\quad\hbox{on }E.
 \tag{HLM.11}
\]
Moreover \(E_0\) is contained in the sum of those original
parts related to a character with the same values as
\(\chi_{x_0}\).
\begin{proof}
Any finite linear relation among the \(s(m)\) lies at one
level containing \(z\).  Applying the given action on that
\(E_K\) shows that the same relation holds among
\(s(\alpha(t)m)\).  Thus these values define a unique
linear operator \(T_t\) on \(E\).  The ring laws and the
commutation with \(\pi\) follow on the spanning set \(s(M)\).
Each vector belongs to an \(E_K\), hence is a finite sum
of simultaneous scalar vectors of the stated kinds.

Eigenspaces for distinct characters are linearly independent.
Indeed, in a finite purported relation and for a chosen
character \(\chi_0\), choose for each other character
\(\chi_j\) some \(t_j\) with
\(\chi_0(t_j)\ne\chi_j(t_j)\).  The commuting operator
\(\prod_j(T_{t_j}-\chi_j(t_j))\) kills all other summands
and is a nonzero scalar on the chosen one.  Thus the relation's
chosen summand is zero.  It follows that \(E\) is the
algebraic direct sum of its character spaces.  At least one
component of \(s(z)\) is nonzero.  Take \(p\) to be that
component projection.  Since \(\pi(g)\) commutes with all
\(T_t\), it preserves every character space and commutes
with \(p\).  Grouping the original spanning subspaces by
character proves the last assertion.  If a map on the entire
ambient \(V\) is desired, extend \(p\) by zero on any
linear complement of \(E\); every asserted identity concerns
\(E\), so this extension introduces no equivariance claim
outside it.
\end{proof}
% END HLM11

% BEGIN HLM12
\paragraph{HLM12. Finite local isotypic refinement preserves all bridges.}
\label{claim:HLM12}
Let \(Q\) be finite.  If a representation of \(\prod_qG_q\)
is, for each factor, a sum of copies of specified irreducible
representations, then every nonzero vector survives a group
equivariant linear map to a representation isotypic for one
specified irreducible at each factor.  Suppose the specified
irreducibles arise from original parts having a common full
determinant law.  One may use the Galois representation of any
one such part in all the resulting local \(\mathrm{WOT}\)
assertions, for every prolongation and every original
\(\mathrm{IsRecT}\) map.
\begin{proof}
A sum of simple modules is semisimple: select a maximal direct
sum of simple submodules; a simple module not contained in it
would have zero intersection with it and enlarge the sum.
Every submodule of a semisimple module is semisimple, as follows
by the same maximal-sum argument, or by selecting complements
to its intersections with finite sums of simple modules.
Group the simple summands by isomorphism class.  A homomorphism
between nonisomorphic simple modules is zero, so their
isotypic sums form a canonical direct decomposition.  Every
endomorphism commuting with this factor preserves that
decomposition.  Consequently its component projections commute
with all the other factors.  Project onto a component on which
the chosen vector is nonzero, and repeat for the finitely many
factors.  Previous isotypic properties persist under taking
submodules and images.  This proves the first assertion,
including \(Q=\varnothing\), when no projection is needed.

For the second assertion, the finite-dimensional Galois
representations of the parts are semisimple.  Equality of their
determinants gives equality of their characters.  Here the
coefficient field has characteristic zero.  To recall the
character argument, take the finite-dimensional algebra generated
by the simultaneous representation on their direct sum.
Its radical acts trivially on both semisimple modules, and its
semisimple quotient is a product of matrix algebras over the
algebraically closed field.  The traces of the central primitive
idempotents give the multiplicities times the simple dimensions;
they determine the multiplicities.  Equal characters on the
group give equal traces on its linear span, hence on these
idempotents.  Thus the two representations are isomorphic.
Conjugation induces the same isomorphism on the period modules
defining \(\mathrm{WOT}\).  It transports the part's local
comparison for each given prolongation and each given
\(\mathrm{IsRecT}\) map.  Nothing in this transport chooses
one preferred local correspondence or weakens the universal
quantifier over them.
\end{proof}
% END HLM12

% BEGIN HLM13
\paragraph{HLM13. The full original \texttt{piPts\_n\_alt}.}
\label{claim:HLM13}
Under the original hypotheses of HLM5, an admitted
\(\lambda_0\) with \(\mathrm{DenseWt}(\theta,\lambda_0)\),
the supplied \(\mathrm{AutomorphicHeckeSet}(\iota,P)\),
and the levelwise \(\mathrm{IsUniv}\) restriction property
of HLM10, every nonzero \(z\in\mathcal M\) admits
\[
 (x,r,Y,\pi,\mathrm{sp},(\sigma_q)_q),
 \quad x:R\longrightarrow\overline{\mathbb Q}_\ell,
 \quad r:H.G_q\longrightarrow\mathrm{GL}_n(\overline{\mathbb Q}_\ell),
 \tag{HLM.12}
\]
where \(Y\) is a complex vector space, \(\pi\) is a
representation of \(\prod_qG_q\), \(\mathrm{sp}\) is
additive, and each \(\sigma_q\) is the original smooth
irreducible local representation, with all the following properties:
\[
 \begin{gathered}
 \mathrm{sp}(z)\ne0,\qquad r\text{ is continuous},\qquad
 D_0\mathbin{\mathrm{map}}x=\det r,\\
 \mathrm{sp}(\alpha(t)m)=\iota(x(t))\mathrm{sp}(m),\qquad
 \mathrm{sp}(\rho(g)m)=\pi(g)\mathrm{sp}(m),\\
 \operatorname{IsoSum}(\pi|_{G_q},Y,\sigma_q)\quad(q\in Q).
 \end{gathered}
 \tag{HLM.13}
\]
For every \(q\), every valuation subring \(B_q\) over
\(H.wq(q)\), and every \(rc\) satisfying its original
\(\mathrm{IsRecT}\), one has
\(\mathrm{WOT}(\iota,r\circ H.\mathrm{prj},
rc(\operatorname{algNorm}\sigma_q))\).
\begin{proof}
Choose a level containing \(z\), and an \(A\)-basis of
its finite free ordinary module \(N_K\).  Some coefficient
of \(z\) is nonzero.  Density supplies a good point
\(x_0:A\to_O\overline{\mathbb Q}_\ell\) at which that
coefficient remains nonzero.  ASP13 gives one specialization
\(s\) on all ambient functions, whose images of this basis
are linearly independent.  Expansion in the basis proves
\(s(z)\ne0\).  On the whole union \(s\) is additive:
any two vectors lie at one level, where ASP13 asserts
additivity.  The same argument gives its scalar rule.

For each level put \(E_K=\operatorname{span}_{\mathbb C}s(M_K)\).
AFS19 gives the actual Hecke action on \(E_K\), its spanning
by intersections with original parts, and a scalar character
on every nonzero such intersection.  At any level containing
\(z\), its Hecke algebra is nonzero and is in \(\mathcal C_O\)
by OHA11.  Compose its specified map \(R\to T_K\) with
each of these scalar characters.  HEP7, applied to a nonzero
vector in the intersection, identifies the resulting full
determinant with that of the part.  Thus these are exactly
the related characters required in HLM11.

The finite product of the actual sections defines the complex
right-translation representation \(\pi_0\) on the ambient
classical functions.  ASP13's action field, applied after
putting a vector and its translate at a common level, gives
\(s\rho(g)=\pi_0(g)s\).  For a single factor the relevant
operators commute: away operators are finite sums of translations
at disjoint sections, and the ordinary operators are finite
sums on the other factor of \(G_a\times G_p\).
Their order is the original right-translation order.  This
also proves the scalar action commutation if only the
levelwise restriction property, rather than HLM10's stronger
global conclusion, was assumed.

Apply HLM11.  Its surviving character is
\(\iota\circ x\) for an actual map \(x:R\to\overline{\mathbb Q}_\ell\),
and its image lies in the sum of parts whose full determinant
is \(D_0\mathbin{\mathrm{map}}x\).  Their local
\(\operatorname{IsoSum}\) fields imply the same spanning
assertion on this subrepresentation, since submodules of a
semisimple module are semisimple.  HLM12 now refines the
finitely many local factors while retaining a nonzero image
of \(s(z)\).  Compose the two projections with \(s\).
Use the continuous Galois representation of the part originally
selected in HLM11.  HLM12 transports all local comparisons
from the parts selected during refinement.  Scalar and group
equivariance survive both projections, proving every field
in (HLM.13) and every stated local comparison.  In particular,
the conclusion is for each module vector, not just for
eigenvectors already chosen at a finite level.
\end{proof}
% END HLM13

% BEGIN HLM14
\paragraph{HLM14. The complete abstract \texttt{umodAbs\_alt} interface.}
\label{claim:HLM14}
Assume precisely the supplied original \(H.\mathrm{Ax}\),
\(\mathrm{IsWt}\), \(\mathrm{AutomorphicHeckeSet}(\iota,P)\),
admitted dense weight, neatness of all levels, and section
faithfulness of HLM5, together with the original literature
hypothesis and the unique-map and all-level \(\mathrm{IsJ}\)
hypotheses of HLM10.  Then the original union \(\mathcal M\)
has the representation \(\rho\) and action \(\alpha\)
of \texttt{umodAbs\_alt}: they commute, the weight algebra
acts by its original scalars, the \(K\)-fixed module is
exactly \(M_K\) and is finite, every normal compact-level
quotient has the specified projective group-algebra action,
the restrictions are actual \(\mathrm{IsMap}\) actions,
and all nonzero vectors have exactly the points of HLM13.
\begin{proof}
HLM5 constructs the representation, fixed-space identity and
all the quotient actions, with the displayed compatibility
on each quotient coset.  HLM10 constructs the single ring
action and proves its restriction identities.  HLM13 uses
those same maps and the same union, so its scalar and group
equivariance are the required ones.  These conclusions do
not require changing either action after choosing a point.
This lists all output fields of the abstract theorem.  In
particular the abstract statement retains its supplied
automorphic and all-level-ideal hypotheses; HLM24 below
constructs those hypotheses for the actual auxiliary levels.
\end{proof}
% END HLM14

\subsection{Changing the allowed ramification without changing the data}

% BEGIN HLM15
\paragraph{HLM15. The original quotient, residual object and pair pullback.}
\label{claim:HLM15}
Let \(V_0,V\) be original levels and
\(h:\operatorname{globOf}(r,V)\to\operatorname{globOf}(r,V_0)\)
an original global-frame homomorphism.  Assume its stated
compatibility with the absolute Galois quotient maps.  Its
restriction
\(p_0:V.\mathrm{frame}.G_0\to V_0.\mathrm{frame}.G_0\)
is continuous and surjective, and
\[
 p_0(V.\mathrm{pr}(g))=V_0.\mathrm{pr}(g).
 \tag{HLM.14}
\]
The original \(r^{\mathrm{up}}\) has global residual
characters \(\bar\chi_i\circ p_0\), unchanged local
\(\bar\psi_{v,i}\) and torsion lifts, and unchanged
weight object \((A,\kappa)\).  A pair on \(r\) pulls
back to a pair on \(r^{\mathrm{up}}\) by
\((D,\psi)\mapsto(D\circ p_0,\psi)\).
It also gives the original \(r_N\), compatibility record,
and pair on \(r_N\) used by \(\mathrm{LevMod}\).
\begin{proof}
Continuity is restriction of \(h.\mathrm{cont}\).  Lift a
point in the target kernel of the quadratic character through
the surjection \(h.p\).  Compatibility of quadratic
characters says that its lift is in the source kernel; this
proves surjectivity of \(p_0\).  The assumed compatibility
of quotient maps gives (HLM.14).

The residual characters remain pairwise distinct because
\(p_0\) is surjective, and their kernels are open by
continuity.  The star involution on group algebras commutes
with the map induced by \(p_0\): on a group basis element
this is exactly preservation of conjugation and of the
cyclotomic character by \(h\), and linearity proves it
on all elements.  Thus residual polarization pulls back.
The local groups, their inertia subgroups, residual characters
and torsion lifts have not changed.  Consequently their
\(\mathrm{IsPsi}\) and inertial-extension conditions are
identical, and the same universal \((A,\kappa)\) represents
the weight functor.

Continuity, reduction, polarization and the ordered ordinary
relations of a pair pull back along these same group-algebra
maps.  For a type place, an element of the new \(I_p(u)\)
comes from the same absolute inertia element and lies in
every \(\ker(\bar\chi_i\circ p_0)\).  Its image under
\(p_0\) therefore belongs to the old \(I_p(u)\).
The old trace relation gives the new one.  This proves the
pair pullback, including all scalar algebras in its polynomial
law fields.  The residual object \(r_N\) in the patching
frame has the same global characters and local data.  Its
type inertia subgroup is contained in that of
\(r^{\mathrm{up}}\), directly from its definition by
the restricted inertia domain.  Restricting that one relation
therefore makes the same \(D,\psi\) a pair on \(r_N\).
These are precisely \texttt{rUp}, \texttt{lamUp},
\texttt{pairUp}, \texttt{rN\_compat} and \texttt{pairN};
no new residual representation or new weight ring is chosen.
\end{proof}
% END HLM15

% BEGIN HLM16
\paragraph{HLM16. The identical type works for every lifted Teichmuller family.}
\label{claim:HLM16}
Under HLM15, the given \(\mathrm{TypeForAlt}(r,\iota,\tau_0)\)
implies \(\mathrm{TypeForAlt}(r^{\mathrm{up}},\iota,\tau_0)\).
If \(V_0.S'\subset V.S'\), let \(j=\mathrm{opUp}\)
keep all ordinary indices and include every surviving away
index.  Then the same \(\theta\), the residual eigenvalues
\(a'=a\circ j\), and the same \(\tau_0\) satisfy the
original new-level \(\mathrm{IsTheta}\),
\(\mathrm{IsWt}\), \(\mathrm{IsRes}\), and
\(\mathrm{UnitaryHeckeData}\) conditions.
\begin{proof}
Write \(k\) for the finite residue field of \(O\) and
\(N=|k|-1\).  Two units \(u,v\) with equal reduction
and \(u^N=v^N=1\) are equal.  Indeed
\[
 (u-v)\sum_{a=0}^{N-1}u^av^{N-1-a}=u^N-v^N=0,
\]
and the sum reduces to \(N\bar v^{N-1}\ne0\), so is
a unit in the local ring.  Apply this with \(v=1\) to
an arbitrary original Teichmuller family \(t'_i\) lifting
\(\bar\chi_i\circ p_0\).  Each \(t'_i\) kills
\(\ker p_0\), hence factors uniquely as \(t_i\circ p_0\).
Surjectivity verifies that \(t_i\) is a Teichmuller lift
of the old character.  Equation (HLM.14) identifies the
restrictions of these two families to the original absolute
inertia at every type place.  Apply the old
\(\mathrm{TypeForAlt}\) to this descended family.  Its
assertion about every original \(\mathrm{IsRecT}\) map
has exactly the same local characters and representation
\(\tau_0\), so is the required new assertion.  The argument
began with an arbitrary new family and retains every such
family and every local bridge.

The torus, inertia and reciprocity maps at \(\ell\) are
unchanged.  The \(\mathrm{IsTheta}\) equation and
\(\mathrm{ContChar}\) condition therefore remain literal
identities.  NTL17 supplies neatness at each new level.
In a stabilizer equation neatness forces both acting elements
to be one, and then both coefficient characters are one;
this is exactly the freeness field of \(\mathrm{IsWt}\).
For \(\mathrm{IsRes.frob}\), the new residual determinant
is the old one composed with \(p_0\), and (HLM.14)
identifies its Frobenius argument.  The generator is indexed
by \(j\), so its prescribed coefficient is \(a\circ j\).
The residual uniformizer formula is unchanged because the
local characters, uniformizer and ordinary indices are
unchanged.  All fields of the new unitary data follow from
these checks.  This proves the literal \texttt{TypeForUp}
and the data part of \texttt{step1\_alt} without requiring
a new integral type.
\end{proof}
% END HLM16

% BEGIN HLM17
\paragraph{HLM17. Exact level equality and faithful auxiliary sections.}
\label{claim:HLM17}
For the original unitary groups, membership in
\(\operatorname{lev}K\) is equivalent to the original
maximal-compact conditions away from \(V.S'\) and the
conditions \(\operatorname{cmp}_{w_q}(u)\in K_q\).
If \(V_0.S'=s.S\) and every new \(K_q\) is maximal
compact, then this subgroup equals the original
\(\operatorname{lev}K_0\), for every original \(K_0\).
The section product also satisfies the full faithfulness
hypothesis in HLM14, for every compact auxiliary level.
\begin{proof}
The forward membership assertion holds for every generator
of \(U_Q\vee\bigvee_q\mathrm{sec}_{w_q}(K_q)\): the
sections have identity component away from their underlying
place and the indicated component at their own place.
For the converse, multiply \(u\) on the left successively
by
\(\mathrm{sec}_{w_q}(\operatorname{cmp}_{w_q}u)^{-1}\).
Each factor is in the proposed level.  It makes the selected
component one and leaves every other selected component,
and all components outside \(V.S'\), unchanged.  After the
finite list is exhausted, the remainder belongs to \(U_Q\)
by its defining conditions.  Reversing the multiplications
proves membership.  This argument uses the actual sections,
not an abstract isomorphic product.

The old auxiliary index set is empty: its places would have
underlying place both in \(V_0.S'=s.S\) and outside
\(s.S\).  At each new selected place the maximal-compact
condition also imposes it at its conjugate.  Indeed the unitary
identity identifies that component with conjugated inverse
transpose, which preserves the integral invertible matrices;
this is the opposite-place check of LTP2.  Places above
\(\ell\) have trivial away component.  Thus the displayed
membership criterion is exactly the old one outside \(s.S\),
in both directions.  This proves literal subgroup equality.

Finally let \(j:\prod_qG_q\to U_b\) be any homomorphism
having the stated section on each individual factor.  Since
the product is finite, it is the product of those sections.
Their support and right-inverse identities give
\(\operatorname{cmp}_{w_q}(j(g))=g_q\).  Applying the
forward membership criterion to \(j(g)\in\operatorname{lev}K\)
therefore yields \(g_q\in K_q\) for every \(q\).
This proves the exact \texttt{levelUR2Faith} quantifier over
all such \(j\).  All assertions also hold for an empty
index set, using the empty product and no deletion steps.
\end{proof}
% END HLM17

\subsection{The exact old-level identification}

% BEGIN HLM18
\paragraph{HLM18. Every ordinary Hecke character occurs in an original part.}
\label{claim:HLM18}
Let \(N=N_K\) be the full ordinary module and
\(T=\operatorname{heckeAlg}(N,H.\mathrm{op})\), before
selecting a residual part.  Let \(x:A\to_O\overline{\mathbb Q}_\ell\)
be an arithmetic specialization of admitted weight, and
\(z:T\to\overline{\mathbb Q}_\ell\) a ring map
restricting to \(x\).  Under the original supplied
automorphic decomposition there is an original part \(Y\)
with
\[
  Y.a(w,j+1)=\iota\bigl(z(T_{w,j})\bigr)
       \quad\hbox{for every away index }(w,j).
 \tag{HLM.15}
\]
\begin{proof}
The algebra \(T\) is finite over \(A\), and \(N\) is
a finite faithful \(T\)-module.  Its support is all of
\(\operatorname{Spec}T\): a finite module localized to
zero is killed by some element outside that prime, which
would put the prime outside the zero set of its annihilator.
Thus the point defined by \(z\) has nonzero fiber by
Nakayama.  Equivalently its maximal-ideal block in
\(N\otimes_{A,x}\overline{\mathbb Q}_\ell\) is nonzero.
In a nonzero module over a local finite-dimensional algebra,
take the last nonzero power of the maximal ideal acting on
the module.  It is killed by that maximal ideal.  A nonzero
vector there is an eigenvector for every \(t\in T\)
with eigenvalue \(z(t)\).  This is the faithful-fiber
argument of AFS9, and does not require flatness of the
coefficient extension.

ASP13 identifies this base-changed module with the complex
span of its specialized basis, so the vector remains nonzero
after applying \(\iota\).  Its operator equations are
exactly the specialized original equations.  CFD12's
classical ordinary decomposition expresses it as a finite
sum of vectors in original parts.  All away operators act
on each part by the stated \(Y.a(w,j+1)\).  Group those
finitely many summands by their away scalar characters.
The character-independence argument in HLM11 shows that
at least one occurring character equals the eigenvector's
character: otherwise a product of separating operators
would annihilate every summand but act as a nonzero scalar
on the eigenvector.  This proves (HLM.15) simultaneously
at every away index.  It is the full
\texttt{levelUR2OrdPt\_alt} assertion; it does not first
assume the character is a point of the prescribed residual
Hecke block.
\end{proof}
% END HLM18

% BEGIN HLM19
\paragraph{HLM19. A dense retained Frobenius set determines the residual part.}
\label{claim:HLM19}
In HLM18 suppose a collection of retained Hecke indices
contains every ordinary index, and its away places have
dense inverse-Frobenius image in the original global group.
Suppose \(a\) has the original residual Frobenius
polynomials.  Then selecting the residual part using only
the retained indices gives exactly the full residual part.
This holds over the original \(A\) with a dense admitted
weight and does not assume a flat specialization.
\begin{proof}
Let \(T_0\subset T\) be the finite algebra generated by
the retained operators.  OHA5--OHA6 give idempotents
\(e_0\in T_0\) and \(e\in T\) selecting the indicated
residual characters.  Full residual conditions imply the
retained conditions, so \(eN\subset e_0N\) and
\(e_0e=e\), by faithfulness.  If equality fails, the
idempotent \(d=e_0(1-e)\) has \(dN\ne0\).  This is
a finite projective, hence finite free, \(A\)-module of
positive rank.  Choose a good point \(x\) at the admitted
weight, which exists by density applied to \(1\).
Then \(dN\otimes_{A,x}\overline{\mathbb Q}_\ell\ne0\).
An eigencharacter of its finite commuting operator algebra
gives \(z:T\to\overline{\mathbb Q}_\ell\) restricting
to \(x\), with \(z(e_0)=1\) and \(z(e)=0\).

Every element of \(T\) is integral over \(A\).  ASP2
gives \(|x(A)|\leq1\), so a monic integral equation
shows \(|z(T)|\leq1\): a root of norm greater than one
would make its top power strictly dominate all lower terms.
The finite set of generators of \(\mathfrak m_A\) has
images bounded by some \(b<1\), hence
\(|z(\mathfrak m_AT)|\leq b\).  In the \(e_0\)-factor,
each retained \(T_i-a_i\) has nilpotent reduction modulo
\(\mathfrak m_A\).  Some power times \(e_0\) thus
lies in \(\mathfrak m_AT_0\), and
\[
                 |z(T_i)-a_i|<1
                 \quad(i\text{ retained}).
 \tag{HLM.16}
\]

HLM18 supplies a part with all the away eigenvalues of
\(z\).  Its continuous Galois representation has the
original Frobenius polynomials.  Choose lifts in \(O\)
of the finitely many values of the residual characteristic
coefficient functions; the resulting functions on the
global group are locally constant.  At every retained
inverse Frobenius, (HLM.16) says that each characteristic
coefficient of the part differs from its prescribed lift
by an element of norm less than one.  The set on which
each such inequality holds is closed as well as open:
the coefficient functions are continuous and the open
unit ball of \(\overline{\mathbb Q}_\ell\) is also
closed.  The dense retained set therefore gives the same
inequalities at every group element.  Apply them at every
old away Frobenius.  In the original Hecke polynomial the
coefficient indexed by \(j+1\) is
\((-1)^{j+1}q^{(j+1)j/2}z(T_{w,j})\).
Here \(|q|_\ell=1\); cancellation gives (HLM.16) at
every old away index as well.  The ordinary indices were
retained, so the inequality now holds at every generator.

Reduce the bounded map \(z\) in the residue field of
\(\overline{\mathbb Q}_\ell\).  Its values on all
generators are exactly the prescribed residual scalars,
and its values on \(A\) are the prescribed residue map.
Since these generate \(T\), its kernel is precisely the
maximal ideal selected by \(e\).  The corresponding
idempotent must map to one, contrary to \(z(e)=0\).
Thus \(d=0\), and the two modules coincide.  The use of
a positive-rank direct summand before specialization is
what ensures that no unsupported faithfulness assertion
after base change enters this argument.
\end{proof}
% END HLM19

% BEGIN HLM20
\paragraph{HLM20. Literal equality of the old module and maximal-compact invariants.}
\label{claim:HLM20}
For the data of HLM15--HLM17, let \(K^{\max}\) denote
the new product of maximal compact groups.  Then every
surviving operator is the same original endomorphism,
\[
 H'.\mathrm{op}_A(i)=H.\mathrm{op}_A(j(i)),\qquad
 M'_{K^{\max}}=M_{K_0}
           \quad\hbox{inside the original function space}.
 \tag{HLM.17}
\]
Here \(H'\) denotes the new Hecke setting, not the
literature-hypothesis record.
\begin{proof}
An away index retains the same split place, section,
uniformizer, matrix and finite-coset sum.  An ordinary
index retains the same local group and positive element.
Thus the operator identity is literal, with \(j\) only
changing the membership proof of an index.  HLM17 and
the unchanged \(\theta,\tau_0\) identify the full
ordinary modules.  Selecting the new residual part on
that module amounts to selecting only the old operators
indexed by \(j\).

The new inverse-Frobenius image is dense by FNL16.
Its image under the continuous surjection \(p_0\) is
dense in the old group.  By (HLM.14) these are exactly
the retained old Frobenius arguments; their residual
polynomials are the old ones by HLM16.  The admitted
dense weights and actual parts are supplied by AFS5
and UPA17 for the old data.  HLM19 consequently
identifies the selected part with the full old part.
This proves (HLM.17), including the zero-module case.
It supplies both statements of the original
\texttt{step2\_alt}, rather than an unspecified module
isomorphism at maximal compact level.
\end{proof}
% END HLM20

% BEGIN HLM21
\paragraph{HLM21. Exact comparison of the determinant pairs.}
\label{claim:HLM21}
Suppose \(T',T\in\mathcal C_O\) are the new and old
nonzero Hecke algebras, with their \(\mathrm{IsUniv}\)
pairs \((D',\psi')\), \((D,\psi)\), and
\(\Phi:T'\to_A T\) preserves every surviving generator.
Then
\[
          D'\mathbin{\mathrm{map}}\Phi=D\circ p_0,
       \qquad \Phi_*\psi'_{v,i}=\psi_{v,i}.
 \tag{HLM.18}
\]
\begin{proof}
An \(O\)-map of these local coefficient rings preserves
the residue map, and therefore sends every power of the
maximal ideal into the corresponding power.  Thus the
pushed determinant and characters remain adically
continuous.  At every new away Frobenius, the two
determinant polynomials in (HLM.18) agree by the generator
identity and (HLM.14), with the same \(q\), coefficient
index and inverse Frobenius.  Density, comparison modulo
each ideal power, and separatedness give equality of
all characteristic coefficients on the whole new group.
HLM6 gives equality as full determinant laws.

The two local characters have the same unit equations
because \(\Phi\) preserves \(A\) and \(\theta\).
Their prefix products at the original uniformizer agree
because all ordinary indices are retained.  Cancel
successive prefix products and the common cyclotomic
factor to get equality of the individual values.
Units and the uniformizer give equality on the local
multiplicative group.  Reciprocity, Weil density and
adic separatedness give equality on the whole
decomposition group.  This is the full original
\texttt{levelUR2BaseEq\_alt}, with its original maps.
\end{proof}
% END HLM21

% BEGIN HLM22
\paragraph{HLM22. The old-level identification respects the universal-ring map.}
\label{claim:HLM22}
Let \(B\) be the given old universal ring, \(B'\) the
new one, and \(\beta:B'\to_OB\) carry its pair to
the pullback of the old pair along \(p_0\), with the
same local characters.  The identity of function modules
in (HLM.17) intertwines the new universal action with
the old action through \(\beta\).  Therefore it gives
the required \(\beta\)-equivariant additive equivalence
between the new maximal-compact fixed space and the old
base module.
\begin{proof}
If the common module is zero the conclusion follows
directly.  Otherwise its two Hecke algebras are in
\(\mathcal C_O\).  Restricting the new generated
algebra of endomorphisms gives an \(A\)-algebra map
\(\Phi:T'\to T\); (HLM.17) identifies each generator
with the corresponding old one.  It need not be
assumed surjective.  Let \(f':B'\to T'\) and
\(f_0:B\to T\) be the two specified maps.  HLM21
shows that \(\Phi f'\) and \(f_0\beta\) take the
new universal pair to exactly the same determinant
and local characters.  The unique-map property of
\(B'\) therefore gives
\(\Phi f'=f_0\beta\).  By definition of \(\Phi\),
this equality is equality of the two actions on the
same functions.  HLM5 identifies the maximal-compact
fixed space with that function submodule, and HLM10
identifies its action with \(f'\).  Their composite
is the asserted additive equivalence and satisfies
the exact scalar equation of \(\mathrm{LevMod.eM\_smul}\).
This proves \texttt{levelUR2SmulKey\_alt} without
changing either ideal or coefficient action.
\end{proof}
% END HLM22

\subsection{One base and the complete all-level package}

% BEGIN HLM23
\paragraph{HLM23. Choose the original Taylor--Wiles base before every auxiliary input.}
\label{claim:HLM23}
Fix the given old \(\mathrm{IsJ}\), its ideal \(J\subset B\),
and its level \(K_0\), with \(V_0.S'=s.S\).  Put
\(D=n[L^+:\mathbb Q]\) and
\(S=O[[\mathrm{Empty}\sqcup\mathrm{Fin}(D)]]\).
There is an original
\(b:\mathrm{TWBase}(\mathrm{pairOf}(B.u),J,\mathrm{Empty},D)\)
whose module is the same old \(M_{K_0}\), with its exact
universal action.  This choice depends only on the fixed
old data, and is made before \(V,h,q,v,w,B_q\).
\begin{proof}
The original full local reciprocity input and KX3 give
coordinates on each local pro-\(\ell\) abelianization,
with inertia the \(d_v\) ramified coordinates and
\(\sum_vd_v=[L^+:\mathbb Q]\).  KX4 fixes extensions
\(\psi^0_{v,i}\) with the prescribed residual and torsion
values.  KX5 identifies the given weight ring with the
power series ring on the triples \((v,i,j)\),
\(0\leq i<n\), \(1\leq j\leq d_v\).  Choose an
inertia lift \(\gamma_{v,j}\) of each coordinate basis
vector.  Enumerate the \(D\) triples once, writing
\(v(k),i(k),\gamma_k\), and put
\(c_k=\psi^0_{v(k),i(k)}(\gamma_k)^{-1}\in O^\times\).
The universal character coordinate formulas give an
\(O\)-algebra isomorphism \(e:S\simeq A\) with
\[
 e(X_k)=c_k\kappa_{v(k),i(k)}(\gamma_k)-1,
 \qquad
 \bar c_k\bar\psi_{v(k),i(k)}(\gamma_k)=1.
 \tag{HLM.19}
\]
This is the source's \texttt{weight\_variables\_exist}
for arbitrary rank \(n\); the number of variables is
not specialized to four.

The given \(\mathrm{IsJ}\) supplies
\(\mu:B\to\operatorname{End}_A(M_{K_0})\).
In its nonzero branch it is the given surjective Hecke
map followed by the faithful inclusion into endomorphisms,
so \(\ker\mu=J\).  In its zero branch it is the unique
zero-module action and \(J=B\), giving the same equality.
Weight compatibility gives
\(\mu(\Lambda.\mathrm{toPair}(B.u)(a))m=am\).
Transport the \(A\)-action along \(e\) and put
\(\lambda_b=\Lambda.\mathrm{toPair}(B.u)\circ e\).
Then (HLM.19) and the inertia restriction identity give
\[
 \lambda_b(X_k)=c_kB.u.\psi_{v(k),i(k)}(\gamma_k)-1.
 \tag{HLM.20}
\]
The resulting \(S\)- and \(B\)-actions have the
required scalar equality.  OHA8 gives finite freeness
over \(A\), hence over \(S\) by transporting a basis.
Finally
\(\operatorname{Ann}_B(M_{K_0})=\ker\mu=J\): an
element is in either side exactly when it acts as zero
on every vector.  Thus for every prime ideal \(\mathfrak p\)
the source's support equivalence is literally
\(J\subset\mathfrak p\Longleftrightarrow J\subset\mathfrak p\).
These verifications give every \(\mathrm{TWBase}\) field,
including zero rank for a zero module.  They are the
general-rank construction underlying JSP8, with the
same ideal rather than a radical or enlargement of it.
\end{proof}
% END HLM23

% BEGIN HLM24
\paragraph{HLM24. Construct the actual package for arbitrary auxiliary data.}
\label{claim:HLM24}
Keep \(b\) fixed as in HLM23.  Now take any of the
\(V,h,q,v,w,B_q\) allowed in the literal
\(\mathrm{LevelAlt}\) quantifiers: \(q\in\mathbb N\),
the \(v_j\) are distinct and outside \(s.S\),
\(V.S'=s.S\cup\operatorname{range}v\), the indicated
valuation subrings lie over both \(w_j\) and \(v_j\),
and all the stated quotient-map and split-decomposition
conditions hold.  There are an original coefficient ring
\(R=B'\in\mathcal C_O\), pair \(p\) on the prescribed
\(r_N\), map \(\beta:B'\to_OB\), original union
\(\mathcal M\), commuting actions \(\alpha,\rho\),
and \(\lambda:A\to_OB'\), with the following properties:
all compact-level fixed modules are finite, every normal
compact-level quotient has the projectivity of HLM14,
the maximal-compact fixed module is \(b.M\) through the
\(\beta\)-equivariant map of HLM22, and every nonzero
vector has all the points in HLM13.
\begin{proof}
For an integral element of \(L^+\), membership in the
place defined by \(B_j\) can be tested after either
field embedding.  The two given place-over assertions
therefore imply \(\operatorname{under}(w_j)=v_j\).
They give exactly the original \(Q=\mathrm{qUp}(V,w)\),
with index set \(\mathrm{Fin}(q)\), distinct underlying
places, outside \(s.S\), and the stated range equality.
In particular \(V_0.S'\subset V.S'\).  All original
split-decomposition conditions remain among the hypotheses;
the level record's split condition supplies the sections
used above.

HLM15 pulls the actual pair \(B.u\) back to a pair
\(p'\) for \(r^{\mathrm{up}}\) on \(B\).  This
is the required nonemptiness witness for KX8, not an
assumed existence of an automorphic auxiliary point.
The global and local finite-cohomology inputs in the
original \(\mathrm{HypothesesH}'\) supply KX6 at the
actual finite set \(V.S'\); the original local
reciprocity and duality inputs supply its KX3--KX4
local interfaces.  KX7--KX8 therefore construct the
actual universal ordinary pair ring \(B'\) for
\(r^{\mathrm{up}}\).  Apply its universal property
to \(p'\) to obtain \(\beta\), with
\[
 \beta_*B'.u.D=B.u.D\circ p_0,
 \qquad \beta_*B'.u.\psi=B.u.\psi.
 \tag{HLM.21}
\]
The pair \(p=\mathrm{pairN}(B'.u)\) on \(r_N\)
is supplied by HLM15.

HLM16 constructs the actual new unitary data using
exactly \(\tau_0,\theta,a\circ j,A,\kappa\).
Put \(\lambda=\Lambda^{\mathrm{up}}.\mathrm{toPair}(B'.u)\).
Its inertia identity says
\[
       \lambda(\kappa_{v,i}(\gamma))=B'.u.\psi_{v,i}(\gamma)
                    \quad(\gamma\in I_v).
 \tag{HLM.22}
\]
FNL16 supplies its original frame axioms, AFS5 all
dense weights, UPA17 the actual automorphic parts with
their original literature inputs, NTL17 every neat
level, and HLM17 section faithfulness.  HLM9 supplies
an \(\mathrm{IsJ}\) at every one of these levels for
this \(B'\), including the zero cases.  The unique-map
hypothesis of HLM14 is the universal property of \(B'\):
two maps with the same pair correspond to the same pair
and hence coincide.  Apply HLM14 at the admitted zero
weight, using the fixed \(\iota\).  It produces all
the stated union, action, finiteness, projectivity and
point data.  HLM20--HLM22 give the specified
maximal-compact identification, then HLM23 identifies
its target with the already fixed \(b.M\).
This proves every raw output of \texttt{BaseObjectHomExt.s3}
with actual suppliers rather than a new all-level
\(\mathrm{IsJ}\) or automorphic-set assumption.
\end{proof}
% END HLM24

% BEGIN HLM25
\paragraph{HLM25. Every \texttt{LevMod} field and full classicality.}
\label{claim:HLM25}
The data of HLM24 give an original
\(\mathcal A:\mathrm{LevMod}(b,h,w)\) satisfying
\(\mathcal A.\mathrm{ClassicalAlt}\).
\begin{proof}
For a local field \(k\) and subgroup \(C\subset k_0^\times\),
the original \(\operatorname{iwC}(k,n,C)\) consists
of integral invertible matrices with upper triangular
reduction and each diagonal residue in \(C\).  It
contains the first principal congruence subgroup, hence
is open; it is a closed subgroup of the compact maximal
compact group, hence compact.  The diagonal entries of
the product of two invertible upper triangular matrices
multiply coordinatewise.  Thus reduction followed by
the diagonal modulo \(C\) is a homomorphism from the
Iwahori group, whose kernel is this subgroup.  It is
therefore normal in the full Iwahori.  For the original
\(C=\operatorname{ellC}(\ell,k_0^\times)\), consisting
of elements of order prime to \(\ell\), take finite
products to obtain the exact normal inclusion
\(U_{\mathrm{Iw},1}\triangleleft U_{\mathrm{Iw}}\).
Their quotient is finite because an open subgroup of
a compact group has finite index.  This verifies the
level assertions also for \(q=0\).

Give \(\mathcal M\) its \(R=B'\)-module structure
through \(\alpha\).  Commutation in (HLM.10) makes
each \(\rho(g)\) \(R\)-linear, without changing its
underlying action on functions.  Give \(R\) its
\(S\)-algebra structure by \(\lambda\circ e\), and
give \(\mathcal M\) its \(S\)-action through \(e\)
and the old \(A\)-action.  Then
\(\alpha(\lambda(e(a)))m=e(a)m\) proves the scalar
tower identity; the \(O\)-tower follows because both
maps are \(O\)-linear.  Equations (HLM.19) and
(HLM.22) give the exact required coordinate equation
\[
  X_k\longmapsto c_k p.\psi_{v(k),i(k)}(\gamma_k)-1.
 \tag{HLM.23}
\]

All fixed subgroups are unchanged by viewing the same
action as \(R\)-linear.  HLM24 gives finiteness over
\(A\) at \(U_{\mathrm{Iw},1}\); transporting
generators along \(e:S\simeq A\) gives finiteness
over \(S\).  Its quotient representation is projective
over \(A[U_{\mathrm{Iw}}/U_{\mathrm{Iw},1}]\).
The coefficient isomorphism induced by \(e\) identifies
this group algebra with the corresponding \(S\)-group
algebra, so it transports a projective splitting and
the quotient action to the exact projectivity field.
This is a ring isomorphism transport, not descent along
a nonflat map.  Its action on every quotient coset is
the original \(\rho(u)\) by HLM5.  The map \(e_M\)
is the HLM24 maximal-compact equivalence followed by
the fixed base identification; HLM22 gives its
\(\beta\)-scalar identity.

For clarity the complete non-instance field dictionary is
\[
\begin{array}{c|l}
\text{field}&\text{constructed value or verification}\\ \hline
r_N,\ \mathrm{compat}&\text{HLM15's original residual object and compatibility}\\
R,\ p& B',\ \mathrm{pairN}(B'.u)\text{ from HLM24}\\
\mathrm{segmentWeightExponent}&\text{equation (HLM.23)}\\
\beta,\ \beta\_\mathrm{pair}&\text{equation (HLM.21), with unchanged local characters}\\
P,\rho&\mathcal M\text{ with the same right-translation action, now }R\text{-linear}\\
\mathrm{fin},\ \mathrm{proj}&\text{the preceding fixed-space and group-algebra arguments}\\
e_M,\ e_M\_\mathrm{smul}&\text{the original-function identification of HLM22}.
\end{array}
\tag{HLM.24}
\]
The ring, algebra, module and tower instance fields were
constructed explicitly above.  Finally suppose a vector
maps to zero under every \(\mathrm{LevPtAlt}\) at the
fixed \(\iota\).  If it were nonzero, HLM13 would
give a point with a nonzero image.  Its determinant,
continuity, scalar, group, local-isotypic and every-bridge
fields are exactly those of \(\mathrm{LevPtAlt}\),
because the underlying actions and pair have not changed.
This is a contradiction.  Thus \(\mathrm{ClassicalAlt}\)
holds.  The argument is the full
\texttt{levelUR2ModPack\_alt} in the source file named
\texttt{NonemptyUnitaryHeckeParts}; it does not conflate
that packing theorem with the separate automorphic
nonvanishing theorem used in the ordinary return argument.
\end{proof}
% END HLM25

% BEGIN HLM26
\paragraph{HLM26. The whole literal all-level \texttt{LevelAlt} supplier.}
\label{claim:HLM26}
Assume the original \(\mathrm{HypothesesH}'\).
For every original coefficient prime \(\ell\), CM
number field \(L\), algebraic closure \(\Omega\),
Hida setting \(s\), rank \(n\geq1\), original level
\(V_0\) and \(Q_0\), coefficient object \(c\),
residual object \(r\), fixed
\(\iota:\overline{\mathbb Q}_\ell\simeq\mathbb C\),
universal ring \(B\), weight object \(\Lambda\),
\(\mathrm{UnitaryHeckeData}\ \mathscr R\), and old
auxiliary level \(K_0\), suppose \(V_0.S'=s.S\).
For every ideal \(J\subset B\) with
\(\mathscr R.\mathrm{IsJ}(B,K_0,J)\), the original
\(\mathrm{Pa.LevelAlt}(B,J)\) holds.
\begin{proof}
Choose \(b\) by HLM23, before introducing any auxiliary
data.  For every \(V,h,q,v,w,B_q\) and every list of
proofs satisfying exactly the hypotheses printed in
HLM24, construct its package by HLM24 and its
\(\mathrm{LevMod}\) and \(\mathrm{ClassicalAlt}\)
proof by HLM25.  These are the existential witness
and the universal continuation in the definition of
\(\mathrm{LevelAlt}\), with
\(D=n[L^+:\mathbb Q]\) and no auxiliary variables
in the fixed base.  The coefficient map, pair,
residual object, type, weight and maximal-compact
identification are exactly the ones checked above.
All primes and all ranks in the original statement
remain allowed, including \(\ell=3\); none of the
projectivity arguments divides by a group order or
by \(n!\).

If \(q=0\), the product group is trivial, the
selected place list and all its local tests are empty,
and HLM17--HLM22 identify the sole fixed module with
the old one.  If the old module is zero, HLM23 uses
\(J=B\) and a zero free base, exactly as its
\(\mathrm{IsJ}\) branch requires.  This does not
assert that every auxiliary union is zero; the proof
of HLM25 still separates every vector of any such
union.  Only when a particular union is zero is its
classicality assertion vacuous.  Thus neither edge
case has been removed from the quantifiers.

The conclusion is the complete original
\texttt{level\_of\_isJ\_unitaryHeckeData}, retaining
the local, automorphic and cohomological foundations
in \(\mathrm{HypothesesH}'\) and in its inherited
proofs.  The later conversion of this family to
\(\mathrm{TWAut}\) and \(\mathrm{HasAut}\) is a
separate consumer; no such conclusion is inferred
merely from the present family.
\end{proof}
% END HLM26

\end{document}
